\chapter{Wnioski i podsumowanie}

\section{Podsumowanie wyników}

W niniejszej pracy przedstawiliśmy kompleksowe podejście do algorytmów uczenia ze wzmocnieniem w markowskich procesach decyzyjnych z dyskontowaniem quasi-hiperbolicznym. Główne osiągnięcia obejmują:

\begin{itemize}
\item Sformułowanie modelu MDP z quasi-hiperbolicznym dyskontowaniem
\item Udowodnienie struktury optymalnych polityk jako par $(\mu^\star, \phi_s^\star)$
\item Przedstawienie pięciu kluczowych twierdzeń zapewniających istnienie i zbieżność optymalnych polityk
\item Opracowanie dwóch algorytmów bezmodelowych: oceny polityki i QH Q-Learning
\item Demonstrację działania algorytmów na modelach testowych
\end{itemize}

\section{Główne wyniki teoretyczne}

Praca dostarcza następujących wyników teoretycznych:

\begin{enumerate}
\item \textbf{Twierdzenie o redukcji polityk:} Każdą politykę zależną od historii można zredukować do polityki dwuskładnikowej $(\mu, \phi_s)$.
\item \textbf{Lemat o kontrakcji:} Operator Bellmana dla polityki stacjonarnej jest kontrakcją z faktorem $\beta$.
\item \textbf{Twierdzenie o zbieżności polityk:} Iteracje wartości zbiegają wykładniczo szybko do wartości optymalnej.
\item \textbf{Twierdzenie o istnieniu optymalnej polityki deterministycznej:} Istnieje deterministyczna optymalna para polityk.
\item \textbf{Twierdzenie o zbieżności QH Q-Learning:} Algorytm QH Q-Learning zbiega prawie na pewno do optymalnych polityk przy odpowiednich warunkach na kroki uczenia.
\end{enumerate}

\section{Wkład praktyczny}

Z praktycznego punktu widzenia praca dostarcza:

\begin{itemize}
\item Algorytmy bezmodelowe działające bez znajomości dynamiki środowiska
\item Metody implementacji wykorzystujące aproksymację stochastyczną o dwóch skalach czasowych
\item Demonstrację zastosowania w modelu zarządzania zasobami
\item Narzędzia do modelowania decyzji z uprzedzeniem teraźniejszości
\end{itemize}

\section{Niespójność czasowa}

Kluczowym odkryciem jest demonstracja niespójności czasowej w optymalnych politykach:

\begin{itemize}
\item Dla $\alpha < 1$ optymalna reguła pierwszego kroku $\mu^\star$ różni się od optymalnej polityki stacjonarnej $\phi_s^\star$
\item Decydent zobowiązany (precommitted agent) osiąga lepsze wyniki niż decydent naiwny stale przeplanowujący
\item Zjawisko to ma istotne implikacje dla zrozumienia zachowań ludzkich i projektowania systemów decyzyjnych
\end{itemize}

\section{Ograniczenia}

Przedstawione podejście ma następujące ograniczenia:

\begin{itemize}
\item Ograniczenie do przypadku decydenta zobowiązanego
\item Założenie o skończonych przestrzeniach stanów i akcji
\item Brak rozważenia przypadków częściowo obserwowalnych
\item Ograniczenie do prostych modeli testowych
\end{itemize}

\section{Kierunki dalszych badań}

Przyszłe badania mogą obejmować:

\begin{itemize}
\item \textbf{Rozszerzenie na POMDP:} Częściowo obserwowalne procesy decyzyjne z quasi-hiperbolicznym dyskontowaniem
\item \textbf{Systemy wieloagentowe:} Interakcje strategiczne między agentami z preferencjami czasowo-niespójnymi
\item \textbf{Ciągłe przestrzenie:} Aproksymacja funkcji z wykorzystaniem sieci neuronowych
\item \textbf{Sophisticated agents:} Równowagi subgame perfect dla decydentów świadomych swojej niespójności
\item \textbf{Estymacja parametrów:} Metody uczenia parametrów $\alpha$ i $\beta$ z danych
\item \textbf{Walidacja empiryczna:} Testy z udziałem ludzi w eksperymentach behawioralnych
\end{itemize}

\section{Znaczenie praktyczne}

Wyniki pracy mają istotne znaczenie dla zastosowań w:

\begin{itemize}
\item \textbf{Ekonomia behawioralna:} Modelowanie decyzji oszczędnościowych i konsumpcyjnych
\item \textbf{Zarządzanie zasobami:} Optymalizacja zapasów z uwzględnieniem ludzkiego czynnika
\item \textbf{Systemy rekomendacyjne:} Balansowanie natychmiastowej satysfakcji z długoterminową wartością
\item \textbf{Finanse:} Planowanie inwestycji i emerytur z uprzedzeniem teraźniejszości
\item \textbf{Opieka zdrowotna:} Zrozumienie przestrzegania długoterminowych terapii
\end{itemize}

\section{Zakończenie}

Niniejsza praca stanowi wkład w teorię i praktykę uczenia ze wzmocnieniem poprzez systematyczne zbadanie markowskich procesów decyzyjnych z quasi-hiperbolicznym dyskontowaniem. Przedstawione algorytmy bezmodelowe wraz z gwarancjami teoretycznymi otwierają nowe możliwości zastosowań w dziedzinach wymagających modelowania realistycznych preferencji czasowych decydentów.

Kluczowa innowacja polega na pokazaniu, że złożone polityki niestacjonarne można sprowadzić do eleganckich struktur dwuskładnikowych, co czyni zarówno analizę teoretyczną, jak i implementację praktyczną wykonalnymi. Demonstracja na modelu zarządzania zasobami potwierdza użyteczność proponowanego podejścia w rzeczywistych zastosowaniach.
